\section{Complete support cohomology of the original maps}


The construction starts with the distinct zeros of \(G(R)\) and retains
both of them. It takes finite matrices with their original coefficients
and declared supports to a diagram of cohomology modules over every
compatible support trajectory. This entire diagram carries a canonical
action of \(G_{\LL}(R)\). Its full-support fibre is exactly the original
kernel with its original minimum metric. Nonzero composites are computed
before a sequence is called a complex.

The foundational input is the original semimodule classification
\cite{secondary}, lines 458--945, especially the source labels
\texttt{thm:fiber-modules}, \texttt{thm:reconstruct-semimodule},
\texttt{thm:equivalence-semimodules-diagrams}, and
\texttt{thm:free-semimodule-boolean-cube}.
All reconstruction assertions needed here are proved below.
The original source and conductor matrices are those of OPG5, CEP2--5,
PCL17--21, and GIP1--6 in the pinned source collection \cite{programme}.
No coordinate support is identified across a basis change without its
explicit connecting maps.

\subsection{The two zeros and cancellation paths}

Let \(R\) be a nonzero commutative ring with identity, and let
\[
 S=G(R)=R\sqcup\{\tau\},\qquad e=0_R\in R\subset S.
 \tag{SCD1}
\]
On the supported copy of \(R\), the operations are the ring operations.
The new element \(\tau\) is the additive identity and multiplicative
absorber. Thus \(r+(-r)=e\), not \(\tau\).
The ring projection \(p:S\to R\) sends \(\tau\) to \(0_R\) and fixes
\(R\). The support character \(\chi:S\to\{0,1\}\) sends \(\tau\) to zero
and every supported scalar to one. Checking whether the factors or
summands are supported proves that both are semiring morphisms.
These are the same maps as in the baseline note \cite{globalization}.

For finite sets \(I,J\), a declared support for an \(R\)-matrix
\(d:R^I\to R^J\) is a relation \(D\subseteq J\times I\) containing
every pair with \(d_{ji}\ne0_R\). It may also contain supported zeros.
Define
\[
 \widetilde d_{ji}=
 \begin{cases}d_{ji}\in R\subset S &(j,i)\in D,\\
 \tau &(j,i)\notin D,
 \end{cases}
 \qquad
 D(A)=\{j:\exists i\in A,\ (j,i)\in D\}.
 \tag{SCD2}
\]
The sparse declaration contains exactly the nonzero entries;
the fully entrywise declaration has \(D=J\times I\).
All results below retain whichever support is declared.

A vector of \(S^I\) is a pair \((A,x)\) with \(A\subseteq I\),
\(x\in R^A\). Zero entries of \(x\) are still supported; outside
\(A\) its entries are \(\tau\). Direct matrix multiplication gives
\[
 \widetilde d(A,x)=(D(A),d_{D(A),A}x).
 \tag{SCD3}
\]
Every row with a declared incident edge receives a supported sum,
including a sum that cancels to \(e\).

\begin{proposition}
For \(d:R^I\to R^J\) and \(f:R^J\to R^K\) with supports \(D,F\),
\[
 p(\widetilde f\widetilde d)=fd,\qquad
 \chi(\widetilde f\widetilde d)_{ki}
 =\bigvee_{j\in J}\bigl(\mathbf1_F(k,j)\wedge\mathbf1_D(j,i)\bigr).
 \tag{SCD4}
\]
If \(fd=0_R\), the lifted composite has entry \(e\) at every
declared length-two path and \(\tau\) elsewhere. It is the global
zero matrix over \(S\) exactly when there are no such paths.
\end{proposition}
\begin{proof}
A product in the sum is supported exactly when both entries are
supported. Its ring value is \(f_{kj}d_{ji}\), even when that product
is zero. A nonempty sum of supported terms is supported and has the
ring sum as its value. Only an empty supported sum is \(\tau\).
This proves every assertion, including rings with zero divisors.
The cancellation certificate \(\sum_j f_{kj}d_{ji}=0_R\) therefore
retains all of its individual paths.
\end{proof}

In particular,
\[
 d_0=\begin{pmatrix}1\\1\end{pmatrix},\qquad
 d_1=\begin{pmatrix}1&-1\end{pmatrix},
 \qquad d_1d_0=0_R,\quad \widetilde d_1\widetilde d_0=(e).
 \tag{SCD5}
\]
The exact morphism \(p\) carries this \(e\) to \(0_R\).
It does not identify \(e\) and \(\tau\) in \(S\).

\subsection{Every compatible support trajectory}

Fix a finite cochain complex with its original based matrices:
\[
 C^n=R^{I_n},\quad d_n:C^n\to C^{n+1},\quad
 d_{n+1}d_n=0_R.
 \tag{SCD6}
\]
Only finitely many \(I_n\) are nonempty. Fix every declared support \(D_n\).
A support trajectory is a sequence \(A=(A_n)_n\), \(A_n\subseteq I_n\),
with \(D_n(A_n)\subseteq A_{n+1}\). Let \(\LL_D\) be their componentwise
ordered set.

\begin{theorem}
The lattice \(\LL_D\) is complete and distributive. Its meet and join
are componentwise intersection and union, its bottom is the empty
trajectory, and its top is \(\mathbf I=(I_n)_n\).
For arbitrary seed sets \(U_n\subseteq I_n\), their smallest containing
trajectory is
\[
 \cl_D(U)_n=U_n\cup
 \bigcup_{m<n}D_{n-1}\cdots D_m(U_m).
 \tag{SCD7}
\]
This closure is monotone, idempotent, and preserves unions.
\end{theorem}
\begin{proof}
An edge starting in a union starts in one member, whose next component
contains its endpoint. An edge starting in an intersection has its
endpoint in every next component. Thus arbitrary unions and intersections
are trajectories, with the stated empty conventions. Set-theoretic
distributivity proves the lattice assertion. Every path endpoint in SCD7
is forced in any containing trajectory. Appending one edge proves that
SCD7 itself is a trajectory. This proves minimality and all closure laws.
\end{proof}

Equivalently the graded vertices \((n,i)\) carry directed edges
\((n,i)\to(n+1,j)\) from \(D_n\). Degrees strictly increase, so
reachability defines a finite partial order. The trajectories are
exactly its upward-closed subsets. Every path remains visible.

\begin{theorem}
For \(A\in\LL_D\) put
\[
 C_A^n=R^{A_n},\qquad d_{n,A}=(d_n)_{A_{n+1},A_n}.
 \tag{SCD8}
\]
These are cochain complexes. Extension by zero \(j_{AB}^n:C_A^n\to C_B^n\)
for \(A\subseteq B\) is an injective cochain map. Consequently
\[
 Z_A^n=\ker d_{n,A},\quad B_A^n=\im d_{n-1,A},\quad
 H_A^n=Z_A^n/B_A^n,\qquad
 h_{AB}^n[x]=[j_{AB}^nx]
 \tag{SCD9}
\]
defines a diagram of \(R\)-modules, with exact composition of its maps.
\end{theorem}
\begin{proof}
If \(i\in A_n\) and \(d_{n,ji}\ne0\), then \(j\in A_{n+1}\).
Thus the restricted intermediate sum removes no nonzero term:
\[
 (d_{n+1,A}d_{n,A})_{ki}
 =\sum_{j\in I_{n+1}}d_{n+1,kj}d_{n,ji}=0_R.
\]
Every declared intermediate path is also included in \(A_{n+1}\),
so its cancellation is retained. The same support condition says
that applying \(d_n\) to a vector supported in \(A_n\) produces
a vector in \(R^{A_{n+1}}\). Therefore extension by zero commutes
with the differential. It maps cycles and boundaries to cycles
and boundaries; this proves the quotient map is well defined.
Linearity, identity, and composition follow from zero extension.
\end{proof}

For \(A\subseteq B\), identify the smaller coordinate modules with
their zero-extended copies. The exact kernels and cokernels are
\[
 \ker h_{AB}^n=(Z_A^n\cap B_B^n)/B_A^n,\qquad
 \operatorname{coker}h_{AB}^n=Z_B^n/(Z_A^n+B_B^n).
 \tag{SCD10}
\]
Indeed a smaller cycle becomes zero precisely when it becomes a
larger boundary, and the image is
\((Z_A^n+B_B^n)/B_B^n\).
For SCD5 there are exactly six trajectories: the empty one; the
last vertex; either middle vertex with the last; both middle vertices
with the last; and the full set. The penultimate one has
\(H^1=R(1,1)\), which becomes a boundary in the full exact complex.
The last-vertex fibre has \(H^2=R\). Thus exactness at full support
does not discard the nontrivial cohomology at smaller supports.

\subsection{Reconstruction and the full \(G_{\LL_D}(R)\) base}

Define the disjoint union, retaining its index even when a module is zero,
\[
 \mathcal H_D^n=\bigsqcup_{A\in\LL_D}H_A^n.
 \tag{SCD11}
\]
For elements in two fibres use
\[
 (A,[x])+(B,[y])
 =\bigl(A\cup B,\ h_{A,A\cup B}[x]+h_{B,A\cup B}[y]\bigr).
 \tag{SCD12}
\]
A supported \(r\in R\) acts within its fibre, and \(\tau\) acts by
the bottom-fibre zero. Define \(\mathcal C_D^n,\mathcal Z_D^n,
\mathcal B_D^n\) from \(C_A^n,Z_A^n,B_A^n\) in the same way.

\begin{theorem}
These operations make \(\mathcal H_D^n\) a \(G(R)\)-semimodule.
Its support lattice is exactly \(\LL_D\), its fibres are exactly
\(H_A^n\), and its transports are \(h_{AB}^n\). It is the quotient
of \(\mathcal Z_D^n\) by the congruence
\[
 (A,x)\sim(B,y)
 \quad\Longleftrightarrow\quad A=B\ \hbox{ and }\ x-y\in B_A^n.
 \tag{SCD13}
\]
Thus different support indices are never identified.
\end{theorem}
\begin{proof}
Three summands may all be transported to \(A\cup B\cup C\).
Functoriality of \(h\) makes either parenthesization equal to
the same sum in that module. Commutativity and the bottom identity
follow in the same way. Linearity of transport proves both scalar
distributive laws. In particular \((r+s)x=rx+sx\) occurs in the
same fibre, also when \(r+s=0_R\). The rules involving \(\tau\)
follow from its action. The supported zero acts by
\(e(A,[x])=(A,0)\), and
\((A,0)+(B,0)=(A\cup B,0)\), identifying the support lattice.
Boundaries transport to boundaries, so SCD13 is preserved under
addition and scalars. Its fibre quotient is precisely \(H_A^n\).
\end{proof}

There is an exact distinction between this construction and a
falsely asserted global square-zero lift. Put
\[
 \mathcal D_n(A,x)=(A,d_{n,A}x),\qquad
 \mathbf z_{n,n+2}(A,x)=(A,0).
 \tag{SCD14}
\]
They are \(S\)-linear and
\(\mathcal D_{n+1}\mathcal D_n=\mathbf z_{n,n+2}\).
For nontrivial \(\LL_D\), this is not the global zero map,
whose output has index \(\varnothing\). The usual preimage of
that global zero is not the union of fibre kernels.
Cohomology is taken in the additive category of \(R\)-module
diagrams on the fixed lattice, as explicitly specified in SCD9,
then reconstructed by SCD11--13.

The exact return map to the original free semimodule is
\[
 E_n:\mathcal C_D^n\to S^{I_n},\qquad E_n(A,x)=(A_n,x).
 \tag{SCD15}
\]
It is \(S\)-linear because supports of sums are unions and
zero extensions commute with addition. It can fail to be injective:
two distinct trajectories may have the same \(n\)-th support.
The original lifted matrix is related to it by
\[
 E_{n+1}\mathcal D_n(A,x)
 =T_{D_n(A_n),A_{n+1}}\bigl(\widetilde d_n E_n(A,x)\bigr).
 \tag{SCD16}
\]
Here \(T\) inserts supported zeros at the additional coordinates,
exactly as in the Boolean-cube diagram of \(S^{I_{n+1}}\).
SCD3 and SCD8 prove this identity at every coordinate.

We now include all support scalars. For the bounded distributive
lattice \(\LL=\LL_D\), define
\[
 G_{\LL}(R)
 =\{(r,\mathbf I):r\in R\}\ \cup\
   \{z_A=(0_R,A):A\in\LL\}\ \subset R\times\LL,
 \tag{SCD17}
\]
with coordinatewise operations
\[
 (r,A)+(s,B)=(r+s,A\cup B),\qquad
 (r,A)(s,B)=(rs,A\cap B).
\]
It is a subsemiring of the product of the ring \(R\) and the lattice
semiring \((\LL,\cup,\cap)\): closure follows because a nonzero
ring coordinate requires full support. The zero is \(z_\varnothing\),
the unit is \((1,\mathbf I)\), and
\((0_R,\mathbf I)=z_{\mathbf I}\) is the supported \(e\).

\begin{theorem}[Canonical action of every support scalar]
Every module diagram on \(\LL\), and in particular the cochain,
cycle, boundary and cohomology diagrams above, has the action
\[
 (r,\mathbf I)(A,x)=(A,rx),\qquad
 z_\lambda(A,x)=(\lambda\cap A,0).
 \tag{SCD18}
\]
This defines a \(G_{\LL}(R)\)-semimodule. Restricting along
\[
 G(R)\longrightarrow G_{\LL}(R),\qquad
 \tau\longmapsto z_\varnothing,\quad
 r\longmapsto(r,\mathbf I)
 \tag{SCD19}
\]
recovers exactly the preceding \(G(R)\)-semimodule.
The scalar map is injective whenever \(\varnothing\ne\mathbf I\).
\end{theorem}
\begin{proof}
The two rules agree at the shared element
\((0_R,\mathbf I)=z_{\mathbf I}\): both give \((A,0)\).
A supported scalar preserves the label and acts linearly. For a
zero-support scalar, additivity on vectors is exactly
\[
 \lambda\cap(A\cup B)=(\lambda\cap A)\cup(\lambda\cap B),
\]
because linear transport sends a zero to a zero. Scalar addition
of two such elements uses the other distributive identity,
\((\lambda\cup\mu)\cap A=(\lambda\cap A)\cup(\mu\cap A)\).
If one scalar is \((r,\mathbf I)\) and the other is \(z_\lambda\),
their sum is \((r,\mathbf I)\); the vector sum is
\((A,rx)+(\lambda\cap A,0)=(A,rx)\), as required.
For two supported scalars, scalar additivity is the ordinary
module law, even if their sum is the supported zero.

For scalar multiplication, two zero-support scalars give
\((\lambda\cap\mu)\cap A\). A supported scalar followed or preceded
by \(z_\lambda\) gives the same zero in \(\lambda\cap A\).
Two supported scalars satisfy the ordinary module multiplication
law, also when their product is zero in \(R\).
The unit acts as identity. The global zero maps to the zero of
the bottom fibre, hence is the required absorber. All semimodule
axioms have now been checked. The homomorphism SCD19 follows
directly from SCD17, and its restriction gives the declared action
of both \(\tau\) and every supported scalar. Injectivity follows
from the distinct support coordinates.
\end{proof}

The maps in SCD14 and the cohomology quotient are linear for this
larger semiring: each takes the zero in a fibre to the zero in the
same fibre. Coordinate extraction SCD15 is equivariant along the
scalar morphism
\[
 G_{\LL}(R)\to G_{\mathcal P(I_n)}(R),\qquad
 (r,\mathbf I)\mapsto(r,I_n),\quad z_A\mapsto z_{A_n},
\]
because the coordinate map \(A\mapsto A_n\) preserves union,
intersection, bottom and top.

\begin{theorem}[Exact additive group completion and arithmetic receiver]
The map
\[
 \epsilon_{\rm full}:\mathcal H_D^n\to H_{\mathbf I}^n,\qquad
 (A,[x])\longmapsto h_{A,\mathbf I}^n[x]
 \tag{SCD20}
\]
is the universal additive group completion. It is also the universal
linear map to an \(R\)-module, for the scalar projection
\(\pi:G_{\LL}(R)\to R,\ (r,A)\mapsto r\).
Thus the original cohomology is recovered exactly on the arithmetic
quotient \(G_{\LL}(R)/(e)\cong R\), while every other fibre remains
in \(\mathcal H_D^n\).
\end{theorem}
\begin{proof}
The displayed map preserves addition and is onto, using the top
fibre. Any additive map from \(\mathcal H_D^n\) to an abelian group
kills every \((A,0)\), since it is additively idempotent and the
target permits cancellation. But
\[
 (A,[x])+(\mathbf I,0)
 =(\mathbf I,h_{A,\mathbf I}[x]).
\]
Applying the map forces its value on \((A,[x])\) to be its value
on that transported top-fibre element. The restriction to the
top fibre is a group homomorphism. This proves unique factorization
through SCD20 for every abelian target, hence the universal
group-completion assertion. For \(R\)-linear targets the same proof
retains scalar linearity.

The scalar projection is a semiring morphism. Its kernel ideal
is the full zero-amplitude fibre \(\{z_A:A\in\LL\}\), which equals
the principal ideal \((e)\): \(ez_A=z_A\), while multiplying \(e\)
by any scalar and adding such products always has zero amplitude.
Moreover addition by \(e=z_{\mathbf I}\) sends every zero-amplitude
scalar to \(e\). The congruence identifying \(e\) with the global
zero therefore identifies every \(z_A\) with zero, and leaves
exactly the ring operations on \(R\). Conversely \(\pi\) separates
all different ring amplitudes. Thus this quotient is exactly \(R\).
Finally SCD20 sends \(z_\lambda(A,[x])\), a fibre zero, to zero
and sends supported scalars to their usual action, proving its
\(\pi\)-linearity.
\end{proof}

\subsection{Gluing, matrix changes, and the space defined by curvature}

For \(A,B\in\LL_D\), the original coordinate inclusions give
a degreewise exact sequence of complexes
\[
 0\to C_{A\cap B}\xrightarrow{x\mapsto(x,-x)}
 C_A\oplus C_B\xrightarrow{(x,y)\mapsto x+y}C_{A\cup B}\to0 .
 \tag{SCD21}
\]
Surjectivity follows by assigning coordinates in \(A\) to the first
summand and those in \(B\setminus A\) to the second. A pair summing
to zero has equal and opposite entries supported in the intersection,
which proves exactness. SCD8 proves compatibility with the differentials.

If \(z\) is a cycle in the union, choose \(z=x+y\) in the two
summands. Then \(d_Ax=-d_By=w\) is an intersection cycle. Set
\[
 \partial[z]=[w],\qquad
 \cdots\to H_{A\cap B}^n\to H_A^n\oplus H_B^n
 \to H_{A\cup B}^n\xrightarrow{\partial}H_{A\cap B}^{n+1}\to\cdots .
 \tag{SCD22}
\]
A new splitting changes \((x,y)\) by \((t,-t)\) with \(t\) in
the intersection, so \(w\) changes by \(dt\). Splitting a primitive
of a boundary shows independence of the representative of \([z]\).
For exactness at the middle term, if a pair of cycles sums to a
boundary, subtract a split primitive to make its sum zero. It then
comes from an intersection cycle. If \(\partial[z]=0\), subtract
an intersection primitive from one splitting term and add it to
the other; both become cycles, proving exactness at the union.
Finally if an intersection cycle maps to boundaries in both larger
spaces, choose primitives with opposite signs. Their sum is a union
cycle mapping to that intersection class. These arguments verify
all three repeating positions and also that each composite is zero.

For an actual cochain map \(F_n:C^n\to C'{}^n\), let \(P_n\)
be a declared support containing its nonzero coefficients, and
let \(D'\) be the target differential supports. Then
\[
 F_\#A=\cl_{D'}((P_n(A_n))_n)
 \tag{SCD23}
\]
preserves joins and bottom, and the same matrices restrict to
cochain maps \(C_A\to C'_{F_\#A}\). This follows from \(d'F=Fd\)
and the defining support closure; hence it induces cohomology maps.
If \(G\) is another map with declared supports \(Q_n\), and the
declared support of \(GF\) is contained in \(Q_nP_n\) while retaining
every nonzero coefficient, then
\[
 (GF)_\#A\subseteq G_\#F_\#A .
 \tag{SCD24}
\]
Indeed the trajectory on the right contains every composite
coordinate before cancellation. Its minimal closed subtrajectory
contains the direct product support. The composite cohomology map
equals the direct one followed by this inclusion, since their
underlying matrices are equal. This is the exact comparison when
a cancellation changes the minimal support. An invertible basis
change gives an isomorphism at full support; it is not asserted
to give a bijection of coordinate support lattices.

The reconstructed map is always \(G(R)\)-linear. Its exact
compatibility with the additional lattice scalars is the transport
of a zero along
\[
 F_\#(\lambda\cap A)\subseteq F_\#\lambda\cap F_\#A.
\]
The left side is the support of the image of \(z_\lambda(A,x)\);
the right side is the support obtained by applying
\(z_{F_\#\lambda}\) to the image of \((A,x)\). Both amplitudes
are zero, and the stated inclusion gives their exact diagram map.
It follows from monotonicity of \(F_\#\). If \(F_\#\) also
preserves intersections and top, it induces a semiring morphism
on the two \(G_{\LL}(R)\) bases, and this comparison is equality.
Thus an arbitrary matrix change has the proved support transport,
without an unproved scalar-base isomorphism.

Now allow arbitrary consecutive matrices with
\[
 \Omega_n=d_{n+1}d_n
 \tag{SCD25}
\]
possibly nonzero, and retain the same trajectory lattice.
For a trajectory \(A\), the restricted composite is exactly
\((\Omega_n)_{A_{n+2},A_n}\), by the intermediate-index proof of
SCD8. Define the forbidden vertices to be those \((n,i)\)
whose \(i\)-th column in \(\Omega_n\) is not the zero vector.
Include all ancestors of these vertices in the forbidden set,
and let \(T\) be its complement.

\begin{theorem}[The curvature defines a maximal exact support space]
The set \(T\) is a trajectory. The trajectories whose restricted
matrices are complexes are exactly
\[
 \{A\in\LL_D:A\subseteq T\}.
 \tag{SCD26}
\]
They form a complete distributive lattice with top \(T\).
\end{theorem}
\begin{proof}
An edge from a non-forbidden vertex cannot reach a forbidden
vertex, since that would make its starting vertex an ancestor
of a forbidden vertex. Hence \(T\) is forward closed.
If column \(i\) of \(\Omega_n\) is nonzero, choose \(k\)
with \((\Omega_n)_{ki}\ne0\). At least one summand
\(d_{n+1,kj}d_{n,ji}\) is nonzero, so there is a declared path
from \((n,i)\) to \((n+2,k)\). Every trajectory containing
the first contains the last, and its restricted curvature
therefore has that nonzero entry. It cannot be a complex.
It also cannot contain an ancestor of the first vertex.
Conversely, a trajectory contained in \(T\) contains no initial
column on which \(\Omega_n\) is nonzero. All restricted composites
are therefore zero. Closure under unions and intersections and
the top assertion follow directly.
\end{proof}

This retains the full original curvature rather than declaring it zero.
There is also an exact correction using submodules beyond coordinate
supports. For consecutive maps \(d,f\),
\[
 d(\ker(fd))=\im d\cap\ker f .
 \tag{SCD27}
\]
Both inclusions follow from \(f(dx)=0\). Restricting the first
domain to \(\ker(fd)\) gives a genuine complex, while retaining
the explicit defect \(fd\) and the exact connecting inclusion.
Thus the obstruction specifies both a coordinate-support lattice
and its potentially larger linear submodule receiver.

\subsection{The original source, conductor and invariant observation}

At every literal original cutoff retain
\[
 N\in\{q-1,q,2q-1,2q\},\qquad
 \Pp_N=\C[S]_{\le N},\qquad E=\C[S]/(\chi_k),
 \quad U_N:\Pp_N\to E,\quad p\mapsto[p].
 \tag{SCD28}
\]
The polynomial \(\chi_k\) is the complete original monic root
polynomial of degree \(q\). The source inner product is the original
one, with its complete mass. Polynomial and primary coordinates
are connected by the actual Vandermonde matrix.

The observation is exactly the original OPG5 map:
\[
 P_k=\frac1{n+1}\sum_{a=0}^{n}(C^{-a})^{\otimes k},\qquad
 L=P_k\Pi^{\otimes k}U^{\otimes k}\iota,\qquad
 B=\im L,\quad\jmath\Lambda=L,\quad K=\ker\Lambda .
 \tag{SCD29}
\]
The \(U\) in this tensor formula is the original period-unit map,
distinct from the remainder \(U_N\). Every original period and
invariant column in \(P_k,\Pi,U,\iota\) is kept. See the exact
\href{https://github.com/KokunoYumeto/zeta-function-research-reader/blob/6cf5b3aa8ff395eba6b5512b28d117334da5767b/workbenches/splitzero-tandem/continuations/20260922-cutoff-profile/RETAINED_COMPLETE_PROOF_SOURCES.tex#L9420}{OPG4--5 definition and original cycle action}.

Let \(Z_{\chi,N}\) be multiplication by the complete \(\chi=\chi_k\),
including its conversion to the original source coordinates.
Use \(\Pp_{N-q}=0\) when \(N=q-1\). Then
\[
 \Pp_{N-q}\xrightarrow{Z_{\chi,N}}\Pp_N
                  \xrightarrow{\Lambda U_N}B
 \tag{SCD30}
\]
is a complex, with degrees zero, one and two. Its cohomology is
\[
 H^0=0,\qquad
 H^1\xrightarrow{\sim}K,\quad[p]\mapsto U_Np,\qquad H^2=0.
 \tag{SCD31}
\]
Multiplication by a monic nonzero polynomial is injective.
Division by \(\chi\) proves
\(\ker U_N=\chi\Pp_{N-q}\); \(U_N\) is onto since \(N\ge q-1\).
A cycle therefore maps to \(K\), every member of \(K\) has a lift,
and two lifts differ by exactly a relation. The map to \(B\)
is onto, proving the last assertion. Applying SCD7--20 to these
complete matrices supplies the full \(G_{\LL_D}(\C)\)-cohomology
object, with precisely \(K\) as its terminal arithmetic fibre.

For the conductor graph, retain its original simple-quartet domain
\(m=1,\ k=4l+1\ge9,\ q=(k+1)^2,\ q'=(k-7)^2\).
No conclusion about a more general multiplicity is imported from
this graph presentation. Retain all coefficients of
\(A=\sum_{a,b=0}^8a_{ab}e_{ab}^{(8)}\), all its shifts \(b_{ab}\),
and the original lower polynomial \(\chi'=\chi_{k-8}\). The prime
in \(\chi'\) denotes that polynomial, not a derivative. The maps are
\[
 (\mathcal T_Ap)(S')=\sum_{a,b=0}^8a_{ab}p(S'+b_{ab}),\qquad
 C_{(i,j),(u,w)}=
 \begin{cases}a_{u-i,w-j}&0\le u-i,w-j\le8,\\
 0&\text{otherwise}.
 \end{cases}
 \tag{SCD32}
\]
The root identity \(\beta'_{ij}+b_{ab}=\beta_{i+a,j+b}\)
gives \(\chi'\mid\mathcal T_A(\chi h)\) for every \(h\).
Thus \(\mathcal T_A\) induces the map on quotient spaces with
matrix \(V_{\beta'}^{-1}CV_\beta\). The transpose \(C^{\mathsf T}\)
is multiplication by the nonzero biform \(A\), hence injective
in the polynomial domain. Therefore \(C\) is onto. Its kernel
\(V_A\subset E\) contains the actual \(K\), by OCF17 and CEP5.

The complete pivot and remaining invariant-image maps are retained:
\[
 CN_k^{\mathsf T}=0,\quad
 S_k^{\mathsf T}N_k^{\mathsf T}=I,\quad
 \im N_k^{\mathsf T}=V_A,\qquad J_R=S_kZ_k .
 \tag{SCD33}
\]
These identities hold in the original primary coordinates.
Here \(Z_k\) contains all additional invariant columns, so after
imposing the conductor equation the complete residual equations
still include \(J_R^{\mathsf T}V_\beta p=0\), as in CG29.
The actual residual observation is
\[
 \Xi=\Lambda|_{V_A}:V_A\to B_A=\Lambda(V_A),\qquad\ker\Xi=K.
\]
In the original fixed residual frame it is \(\Xi_k^{\rm fix}\);
changing to that frame retains its full invertible target map.
No formula below substitutes a historical rank numeral for
\(\dim K\) or \(\dim B_A\).

Let \(v\) be the first nonzero translation moment order. The actual
invertible source map \(S_N\) takes a polynomial to its first \(v\)
original source coordinates and its \(\mathcal T_A\) image.
To see invertibility, Taylor's finite expansion says
\(\mathcal T_A\) lowers each degree at least \(v\) by exactly \(v\)
and kills exactly \(\Pp_{v-1}\). The first \(v\) source coordinates
are an isomorphism on that kernel. The source and target dimensions
are equal, so the displayed injective map is invertible.
Define the original graph maps
\[
 \begin{gathered}
 X_N=\C^v\oplus\C[S']_{\le N-v-q'},\qquad
 D_N=\operatorname{diag}(I_v,M_{\chi'}),\\
 E_N=U_NS_N^{-1}D_N:X_N\to V_A,\qquad
 B_N=\Xi E_N:X_N\to B_A .
 \end{gathered}
 \tag{SCD34}
\]
Here \(D_N\) includes the exact conversion to the lower original
source coordinates. The equality
\(S_N^{-1}D_NX_N=U_N^{-1}(V_A)\) follows from divisibility by
\(\chi'\). Consequently \(E_N\) is onto \(V_A\).

The full original relation columns are
\[
 \begin{gathered}
 \mathcal Z_Nh=
 \begin{pmatrix}P_{<v}Z_{\chi,N}h\\\mathcal U_{A,N-q}h\end{pmatrix},
 \qquad D_N\mathcal Z_N=S_NZ_{\chi,N},\\
 (\mathcal U_{A,N-q}h)(S')
 =\sum_{a,b=0}^8a_{ab}d_{ab}(S')h(S'+b_{ab}),\qquad
 d_{ab}(S')=\chi(S'+b_{ab})/\chi'(S').
 \end{gathered}
 \tag{SCD35}
\]
Each \(d_{ab}\) is the original complete product of root differences.
Multiplication by \(\chi'\) proves the lower equality using SCD32;
the upper block is exactly the retained lower source coordinates.
Thus \(\mathcal Z_N\) is injective and
\(\ker E_N=\im\mathcal Z_N\).

The correct graph complex and its full-support cohomology are
\[
 \Pp_{N-q}\xrightarrow{\mathcal Z_N}X_N\xrightarrow{B_N}B_A,
 \qquad H^1\xrightarrow{\sim}K,\quad[x]\mapsto E_Nx,
 \qquad H^0=H^2=0 .
 \tag{SCD36}
\]
Injectivity, surjectivity and \(\ker\Xi=K\) prove these statements
exactly as in SCD31. The chain map to SCD30 is
\[
 (I,\ S_N^{-1}D_N,\ B_A\hookrightarrow B).
 \tag{SCD37}
\]
Its squares are SCD35 and
\(\Lambda U_NS_N^{-1}D_N=\Xi E_N\). Its induced middle
cohomology map is the identity of \(K\). On support trajectories
the exact map is SCD23 with the actual coefficient supports of
these three maps. Hence the support comparison is proved as well.

Choose the full original kernel frame \(I_K\) with actual rank
\(r=\dim K\). In original orthonormal source coordinates let
\(G_N=(U_NU_N^*)^{-1}\), and set
\[
 W_{K,N}=D_N^{-1}S_NU_N^*G_NI_K,\qquad
 E_NW_{K,N}=I_K .
\]
The inverse of \(D_N\) is used only on its indicated image;
the minimum lift is in \(U_N^{-1}(V_A)\). Then the exact kernel is
\[
 \ker B_N=\im[\mathcal Z_N,W_{K,N}],\qquad
 \rk[\mathcal Z_N,W_{K,N}]=(N-q+1)+r .
 \tag{SCD38}
\]
Applying \(E_N\) to a relation among those columns first kills
the coefficients of the independent \(I_K\). Injectivity of
\(\mathcal Z_N\) then kills the remainder. Conversely subtract
the corresponding \(W_{K,N}\) lift from any vector in \(\ker B_N\);
what remains lies in \(\ker E_N\). This proves every column and
rank assertion, including the empty relation list at \(N=q-1\).

The longer sequence \(X_N\xrightarrow{E_N}V_A\xrightarrow{\Xi}B_A\)
has the exact curvature
\[
 \Xi E_N=B_N:X_N\twoheadrightarrow B_A .
 \tag{SCD39}
\]
It is a complex only when the actual \(B_A\) is zero. In all cases
its zero-curvature submodule is precisely SCD38. Thus the nonzero
defect of that proposed sequence is retained, while SCD36 gives
the genuine original graph complex.

\subsection{Exact return to the original metric and arithmetic action}

The full-support identification preserves the actual minimum metric.
For a fixed \(x\in K\), all representatives in SCD30 differ by
\(\chi h\). Minimizing their complete source norm is precisely
the definition of the original quotient norm on \(E\). Thus in
the unchanged full kernel frame it is
\[
 H^K_N=I_K^*G_NI_K .
 \tag{SCD40}
\]
For the graph complex, its actual source metric is
\[
 H_N^{\rm actual}=D_N^*S_N^{-*}S_N^{-1}D_N .
 \tag{SCD41}
\]
The map \(S_N^{-1}D_N\) in SCD37 is then an isometry onto its
image. Every representative of \(x\in K\) lies in
\(U_N^{-1}(V_A)\), so the same minimum gives exactly SCD40.
This proof does not substitute the alternative graph product
metric \(D_N^*D_N\) for SCD41. For all four cutoffs, the full
cohomology determinant is therefore the original determinant,
including the original signed return
\[
 \begin{split}
 &\log\det H^{\rm coh}_{q-1,\mathbf I}
 +\log\det H^{\rm coh}_{q,\mathbf I}
 -\log\det H^{\rm coh}_{2q-1,\mathbf I}
 -\log\det H^{\rm coh}_{2q,\mathbf I}\\
 &\hspace{8mm}=
 \log\det H^K_{q-1}+\log\det H^K_q
 -\log\det H^K_{2q-1}-\log\det H^K_{2q}.
 \end{split}
 \tag{SCD42}
\]
Use the same \(I_K\) at every endpoint and the original empty
determinant convention. This exact identity assigns no missing
asymptotic coefficient.

The original multiplication \(M:E\to E\) need not preserve \(K\).
Its failure to descend through observation is explicitly computable.
Keep the full minimum section and observed action
\[
 Q_N=(\Lambda G_N^{-1}\Lambda^*)^{-1},\quad
 L_N=G_N^{-1}\Lambda^*Q_N,\quad P_N=L_N\Lambda,\quad
 M_{B,N}=\Lambda ML_N .
\]
Then
\[
 \Lambda M-M_{B,N}\Lambda
 =\Lambda M(I-P_N)
 =\Lambda MI_K(H_N^K)^{-1}I_K^*G_N .
 \tag{SCD43}
\]
For proof, direct multiplication gives the first equality.
The minimum section is orthogonal to \(\ker\Lambda=K\);
hence \(I-P_N\) is the \(G_N\)-orthogonal projector onto \(K\).
The displayed expression \(I_K(H_N^K)^{-1}I_K^*G_N\)
is the same projector: multiply to verify idempotence, image,
identity on \(K\), and weighted self-adjointness. This proves
the final equality without deleting its kernel block.

The exact action-compatible subspace and defect space are
\[
 0\to K\cap M^{-1}K\to K
 \xrightarrow{\Lambda M}\im(\Lambda M|_K)\to0.
 \tag{SCD44}
\]
The kernel assertion is its definition; the last target is the
actual image. Thus this is the exact obstruction to obtaining
an action on the original middle cohomology by \(M\).

Retain every higher observation as well:
\[
 K^{(j)}=\bigcap_{a=0}^{j}\ker(\Lambda M^a),\qquad j\ge0.
 \tag{SCD45}
\]
The connecting maps are inclusions, and
\(K^{(j+1)}=K\cap M^{-1}K^{(j)}\). If two successive spaces
are equal, the common space is \(M\)-invariant: membership in
\(K^{(j+1)}\) implies that the \(M\)-image belongs to \(K^{(j)}\).
It is then killed by every further observation, so the sequence
stabilizes. It has at most \(\dim K\) strict drops. Its stable
value is the largest \(M\)-invariant subspace of \(K\), since
every such subspace is killed by every \(\Lambda M^a\), while
the stable value itself is invariant. This is a specified space
with exact connecting maps, rather than an assumed descent.

The original metrics, periods and roots are unchanged throughout.
The construction provides a complete support object over the
absolute base, its exact arithmetic receiver, and the spaces
defined by the two original curvature defects SCD39 and SCD43.
No support trajectory or additional invariant column is discarded.

\subsection{One support base for the four original cutoffs}

The four support bases need not be identified by choosing bases again.
There is a single exact common construction. For a fixed original source
order \(s\in\{1,k\}\), retain the original source inclusions
\(i_{NM}:\Pp_N\hookrightarrow\Pp_M\), \(N<M\), at the four cutoffs.
On relation polynomials use \(j_{NM}:\Pp_{N-q}\hookrightarrow\Pp_{M-q}\)
and on \(B\) use the identity. Their original matrices satisfy
\[
 i_{NM}Z_{\chi,N}=Z_{\chi,M}j_{NM},\qquad
 U_Mi_{NM}=U_N,\qquad
 \Lambda U_Mi_{NM}=\Lambda U_N .
 \tag{SCD46}
\]
These are identities on polynomials before passing to any coordinates.
Their matrices include whatever original basis conversions are required.
They form cochain maps of SCD30 and induce precisely the identity of
\(K\) in SCD31.

For the graph presentations, the corresponding map is
\[
 F^X_{NM}=D_M^{-1}S_Mi_{NM}S_N^{-1}D_N,\qquad
 E_MF^X_{NM}=E_N,\quad
 F^X_{NM}\mathcal Z_N=\mathcal Z_Mj_{NM}.
 \tag{SCD47}
\]
The inverse of \(D_M\) is taken only on its actual image: a polynomial
whose class lies in \(V_A\) still has that property after source inclusion,
so \(S_M\) of that polynomial has its second component divisible by
\(\chi'\). This proves that the map is defined. Cancellation of the
displayed inverse maps, followed by SCD35 and SCD46, proves the two
identities. Composition for three source cutoffs follows by the same
cancellation. The induced cohomology map is again the identity of \(K\).

Now form a finite coefficient graph whose vertices record the cutoff,
cochain degree and coefficient index. Include every declared differential
edge, every edge of these original horizontal source maps, and, if both
presentations are retained together, every edge of the graph-to-polynomial
map SCD37. Let \(\LL_{\rm all}\) be all subsets closed under these edges.
One can include both source orders as separately labelled copies, without
identifying their metrics. The proof of SCD7 applies to this finite
diagram: unions and intersections remain closed, and forward closure is
the least containing closed set. Thus \(\LL_{\rm all}\) is a finite bounded
distributive lattice with all the original coefficient supports retained.

The coordinate projection to any one of the complexes preserves union,
intersection, bottom and top, and takes a master trajectory to a valid
local trajectory. Restricting every original matrix at a master trajectory
therefore gives a diagram of complexes in which all original squares
still commute. To verify a square entry, note that every intermediate
index with a nonzero contribution is forced into the master trajectory
by its first edge; its second edge is retained as well. Thus the restricted
entry is the original full entry, with all its cancellation terms.

For each cutoff reconstruct
\[
 \mathcal H_{\rm all,N}^1
 =\bigsqcup_{A\in\LL_{\rm all}}H^1_{N,A|_N}.
 \tag{SCD48}
\]
These are modules over the \emph{same} \(G_{\LL_{\rm all}}(\C)\) by
SCD18. The horizontal cohomology maps preserve the master label \(A\),
commute with union transports, and take a zero in a fibre to the zero in
the same fibre. They are therefore \(G_{\LL_{\rm all}}(\C)\)-linear:
for \(z_\lambda\) both composites give the zero at \(\lambda\cap A\);
for a supported scalar the assertion is ordinary linearity.
Their additive group completions all equal the same \(K\), with identity
connecting maps and the four separate original metrics SCD40.
This proves a common absolute-base receiver for the complete four-cutoff
diagram, including the graph maps, without merging its source metrics.


\subsection{Sources and verification scope}

The accompanying exact checker verifies the complete six-trajectory
cancellation example and its cohomology changes; the \(G_{\LL}(R)\)
scalar laws; and the full source/graph/minimum identities on a
specified positive four-root auxiliary example at all four cutoffs.
This example tests the proved algebra. It is not a replacement
for the original period matrices or an evaluation of a hypothetical
zeta zero. No floating-point test proves a native sign or coefficient.

The classification source is read as original LaTeX, not a PDF or
a summary. The accompanying source-use ledger records exact hashes,
reading coverage and receiving formulas. The elementary group-completion,
gluing and cohomology constructions above have complete proofs here;
the human semiring reference below is contextual provenance, not a
claim that an unread external theorem supplies a missing calculation.

